% !TEX spellcheck = en_US
% !TEX spellcheck = LaTeX
\documentclass[letterpaper,10pt,english]{article}
\usepackage{%
amsfonts,%
amsmath,%
amssymb,%
amsthm,%
babel,%
bbm,%
%biblatex,%
caption,%
centernot,%
color,%
enumerate,%
%enumitem,%
epsfig,%
epstopdf,%
etex,%
fancybox,%
framed,%
fullpage,%
%geometry,%
graphicx,%
hyperref,%
latexsym,%
mathptmx,%
mathtools,%
multicol,%
paralist,%
pgf,%
pgfplots,%
pgfplotstable,%
pgfpages,%
proof,%
psfrag,%
%subfigure,%
tcolorbox,%
tfrupee,%
tikz,%
times,%
%ulem,%
url,%
xcolor,%
mathpazo,%
}

\definecolor{shadecolor}{gray}{.95}%{rgb}{1,0,0}
\usepackage[margin=1in,top=0.75in]{geometry}
\usepackage[mathscr]{eucal}
\usepgflibrary{shapes}
\usepgfplotslibrary{fillbetween}
\usetikzlibrary{%
arrows,%
backgrounds,%
chains,%
decorations.pathmorphing,% /pgf/decoration/random steps | erste Graphik
decorations.text,% 
matrix,%
positioning,% wg. " of "
fit,%
patterns,%
petri,%
plotmarks,%
scopes,%
shadows,%
shapes.misc,% wg. rounded rectangle
shapes.arrows,%
shapes.callouts,%
shapes%
}

%\pgfplotsset{compat=newest} %<------ Here
\pgfplotsset{compat=1.11} %<------ Or use this one

\theoremstyle{plain}
\newtheorem{thm}{Theorem}[section]
\newtheorem{lem}[thm]{Lemma}
\newtheorem{prop}[thm]{Proposition}
\newtheorem{cor}[thm]{Corollary}
\newtheorem{clm}[thm]{Claim}

\theoremstyle{definition}
\newtheorem{defn}[thm]{Definition}
\newtheorem{conj}[thm]{Conjecture}
\newtheorem{exmp}[thm]{Example}
\newtheorem{axiom}[thm]{Axiom}
\newtheorem{assum}[thm]{Assumption}
\newtheorem{exerc}[thm]{Exercise}
\newtheorem*{defin}{Definition}

\theoremstyle{remark}
\newtheorem{rem}{Remark}
\newtheorem{note}{Note}

\newcommand{\iid}{\emph{i.i.d.}\ }
\newcommand{\Cov}{\operatorname{Cov}}
%\newcommand{\det}{\operatorname{det}}
%\newcommand{\Real}{\mathbb{R}}
\newcommand{\tr}{\operatorname{tr}}
\newcommand{\Var}{\operatorname{Var}}
\newcommand{\cov}{\operatorname{cov}}

%\renewcommand{\proof}[1]{\begin{proof}#1\end{proof}}
\newcommand{\EQ}[1]{\begin{equation*}#1\end{equation*}}
\newcommand{\EQN}[1]{\begin{equation}#1\end{equation}}
\newcommand{\eq}[1]{\begin{align*}#1\end{align*}}
\newcommand{\meq}[2]{\begin{xalignat*}{#1}#2\end{xalignat*}}
\newcommand{\norm}[1]{\left\lVert#1\right\rVert}
\newcommand{\inner}[1]{\left\langle #1\right\rangle}
\newcommand{\abs}[1]{\left\lvert#1\right\rvert}
\newcommand{\expect}[1]{\mathbb{E}\left[{#1}\right]}
\newcommand{\prob}[1]{\mathbb{P}\left[{#1}\right]}
\newcommand{\given}{\; \big\vert \;} 
\newcommand{\set}[1]{\left\{#1\right\}} 
\newcommand{\indicator}[1]{1_{\set{#1}}} 
\newcommand{\Ind}[1]{\mathbbm{1}_{#1}} 
\newcommand{\SetIn}[1]{\mathbbm{1}_{\set{#1}}} 
\newcommand{\red}[1]{\textcolor{red}{#1}} 
\newcommand{\floor}[1]{\left\lfloor#1\right\rfloor}
\newcommand{\ceil}[1]{\left\lceil#1\right\rceil}


\newcommand{\C}{\mathbb{C}}
\newcommand{\D}{\mathbb{D}}
\newcommand{\E}{\mathbb{E}}
\newcommand{\F}{\mathbb{F}}
\newcommand{\N}{\mathbb{N}}
\renewcommand{\P}{\mathbb{P}}
\newcommand{\Q}{\mathbb{Q}}
\newcommand{\R}{\mathbb{R}}
\newcommand{\Z}{\mathbb{Z}}

\newcommand{\bA}{\mathbf{A}}
\newcommand{\bI}{\mathbf{I}}
\newcommand{\bU}{\mathbf{1}}
\newcommand{\bx}{\mathbf{x}}
\newcommand{\bX}{\mathbf{X}}
\newcommand{\bZ}{\mathbf{Z}}


\newcommand{\cA}{\mathcal{A}}
\newcommand{\cB}{\mathcal{B}}
\newcommand{\cC}{\mathcal{C}}
\newcommand{\cF}{\mathcal{F}}
\newcommand{\cG}{\mathcal{G}}
\newcommand{\cM}{\mathcal{M}}
\newcommand{\cN}{\mathcal{N}}
\newcommand{\cP}{\mathcal{P}}
\newcommand{\cT}{\mathcal{T}}
\newcommand{\cX}{\mathcal{X}}
\newcommand{\cY}{\mathcal{Y}}

\newcommand{\sA}{\mathscr{A}}
\newcommand{\sB}{\mathscr{B}}
\newcommand{\sC}{\mathscr{C}}
\newcommand{\sE}{\mathscr{E}}
\newcommand{\sF}{\mathscr{F}}
\newcommand{\sG}{\mathscr{G}}
\newcommand{\sH}{\mathscr{H}}
\newcommand{\sL}{\mathscr{L}}
\newcommand{\sS}{\mathscr{S}}
\newcommand{\sT}{\mathscr{T}}
\newcommand{\sX}{\mathscr{X}}
\newcommand{\sY}{\mathscr{Y}}
\newcommand{\sZ}{\mathscr{Z}}

\newcommand{\tS}{\tilde{S}}
% Debug
\newcommand{\todo}[1]{\begin{color}{blue}{{\bf~[TODO:~#1]}}\end{color}}

% a few handy macros

\renewcommand{\le}{\leqslant}
\renewcommand{\ge}{\geqslant}
\newcommand\matlab{{\sc matlab}}
\newcommand{\goto}{\rightarrow}
\newcommand{\bigo}{{\mathcal O}}
%\newcommand{\half}{\frac{1}{2}}
%\newcommand\implies{\quad\Longrightarrow\quad}
\newcommand\reals{{{\rm l} \kern -.15em {\rm R} }}
\newcommand\complex{{\raisebox{.043ex}{\rule{0.07em}{1.56ex}} \hskip -.35em {\rm C}}}


% macros for matrices/vectors:

% matrix environment for vectors or matrices where elements are centered
\newenvironment{mat}{\left[\begin{array}{ccccccccccccccc}}{\end{array}\right]}
\newcommand\bcm{\begin{mat}}
\newcommand\ecm{\end{mat}}

% matrix environment for vectors or matrices where elements are right justifvied
\newenvironment{rmat}{\left[\begin{array}{rrrrrrrrrrrrr}}{\end{array}\right]}
\newcommand\brm{\begin{rmat}}
\newcommand\erm{\end{rmat}}

% for left brace and a set of choices
%\newenvironment{choices}{\left\{ \begin{array}{ll}}{\end{array}\right.}
\newcommand\when{&\text{if~}}
\newcommand\otherwise{&\text{otherwise}}
% sample usage:
%\delta_{ij} = \begin{choices} 1 \when i=j, \\ 0 \otherwise \end{choices}


% for labeling and referencing equations:
\newcommand{\eql}{\begin{equation}\label}
\newcommand{\eqn}[1]{(\ref{#1})}
% can then do
%\eql{eqnlabel}
%...
%\end{equation}
% and refer to it as equation \eqn{eqnlabel}.
% some useful macros for finite difference methods:
\newcommand\unp{U^{n+1}}
\newcommand\unm{U^{n-1}}
% for chemical reactions:
\newcommand{\react}[1]{\stackrel{K_{#1}}{\rightarrow}}
\newcommand{\reactb}[2]{\stackrel{K_{#1}}{~\stackrel{\rightleftharpoons}
 {\scriptstyle K_{#2}}}~}
\makeatletter
\def\th@plain{%
\thm@notefont{}% same as heading font
\itshape % body font
}
\def\th@definition{%
\thm@notefont{}% same as heading font
\normalfont % body font
}
\makeatother
\date{}


\title{Lecture-05: Random Vectors}
\author{}

\begin{document}
\maketitle
\section{Random vectors} 
\begin{defn}[Projection] 
For a vector $x \in \R^n$, we can define $\pi_i: \R^n \to \R$ is the \textbf{projection} of an $n$-length vector onto its $i$-th component, such that $\pi_i(x) = x_i$. 
\end{defn}
%\begin{defn}
%For a subset $B \subseteq \R^n$, we define sets $\pi_i(B) \triangleq \set{x \in \R^n: \pi_i(x) = \pi_i(b) \text{ for some } b \in B}$.
%\end{defn}
\begin{defn}
The Borel sigma algebra over space $\R^n$ is defined as the smallest sigma algebra generated by the family $(\pi_i^{-1}(B_x): x \in \R, i \in [n])$ and is denoted by $\sB(\R^n)$. 
The elements of the Borel sigma algebra are called Borel sets. 
\end{defn}
\begin{rem}
By definition of $\sB(\R^n)$, projection $\pi_i: \R^n \to \R$ is a Borel measurable function for all $i \in [n]$. 
For a subset $A \subseteq \R$ and projection $\pi_i: \R^n \to \R$, 
we can write 
\EQ{
\pi_i^{-1}(A) = \set{x \in \R^n: x_i \in A} = \R \times \dots \times A \times \dots \times \R.
}
Thus, for any $A \in \sB(\R)$, we have $\pi_i^{-1}(A) \in \sB(\R^n)$. 
%It follows that the function $\pi_i: \R^n \to \R$ is a Borel measurable function.
\end{rem}
%\begin{rem} 
%For any $A \subseteq \R$, we see that $\pi_i^{-1}(A)= \R \times \dots \times A \times \dots \times \R$. 
%\end{rem}
%\begin{shaded}
%\begin{exmp}
%Consider a function $f: \Omega \to \R^n$, then $f(\omega) \in \R^n$ and can be expressed in terms of its components $(f_1(\omega), \dots, f_n(\omega))$, where $f_i = \pi_i\circ f$. 
%For this function $f$, we can write the inverse image of a set $B \subseteq \R^n$ as 
%\EQ{
%f^{-1}(B) 
%\triangleq \cap_{i=1}^n\set{\omega \in \Omega: (\pi_i\circ f)(\omega) \in \pi_i(B)}
%= \cap_{i=1}^nf_i^{-1}(\pi_i(B)).
%}
%%For the component function $f_i: \Omega \to \R$ and a set $A \subseteq \R$, 
%%we can write 
%%\EQ{
%%f_i^{-1}(A) 
%%\triangleq \set{\omega \in \Omega: f_i(\omega) \in A} 
%%= \set{\omega \in \Omega: (\pi_i \circ f)(\omega) \in A}
%%= \set{\omega \in \Omega: f(\omega) \in \pi_i^{-1}(A)} 
%%= (f^{-1} \circ \pi_i^{-1})(A). 
%%}
%\end{exmp}
%\end{shaded}
%%Consider the probability space $(\Omega, \sF, P)$, 
%%and a random vector $X: \Omega \to \R^n$. 
%%Let $\pi_i: \R^n \to \R$ be the projection of an $n$-length vector to its $i$th component, i.e. $\pi_i(X) = X_i$. 
%%\begin{shaded}
%%\begin{exerc} 
%%Show that the function $\pi_i: \R^n \to \R$ defined by $g: x \mapsto x_i$ is a Borel measurable function. 
%%\end{exerc}
%%\end{shaded}
%
%%\begin{exerc} 
%%Show that $(\pi_i^{-1}\circ \pi_i)(B) = \R \times \dots B \dots \times \R$.
%%Since $f_i^{-1} = f^{-1}\circ \pi_i^{-1}$ and $(\pi_i^{-1}\circ \pi_i)(B) = \R \times \dots \pi_i(B) \dots \times \R$, 
%%we can write $f_i^{-1}(\pi_i(B)) = \R \times \dots \pi_i(f^{-1}(B)) \dots \times \R$. 
%%\end{exerc}
\begin{defn}[Random vectors] 
Consider a probability space $(\Omega, \sF, P)$ and a finite $n \in \N$. 
A \textbf{random vector} $X: \Omega \to \R^n$ is an $\sF$-measurable mapping from the sample space to an $n$-length real-valued vector. 
That is, for any $x \in \R^n$, we have %the event 
\EQ{
A_X(x) \triangleq \set{\omega \in \Omega: X_1(\omega) \le x_1, \dots, X_n(\omega) \le x_n} 
= \cap_{i=1}^nX_i^{-1}(-\infty, x_i] 
\in \sF.
}
%We say that the random vector $X$ is $\sF$-measurable.  
\end{defn}
\begin{shaded}
\begin{exmp}[Tuple of indicators] 
Consider a probability space $(\Omega, \sF, P)$, a finite $n \in \N$, and events $A_1, \dots, A_n \in \sF$. 
We define a mapping $X: \Omega \to \set{0,1}^n$ by $X_i(\omega) \triangleq \Ind{A_i}(\omega)$ for all outcomes $\omega \in \Omega$. 
Let $x \in \R^n$, then we can write $A_X(x) = \cap_{i=1}^n\Ind{A_i}^{-1}(-\infty, x_i]$. 
Recall that 
\EQ{
\Ind{A_i}^{-1}(-\infty, x_i] = 
\begin{cases}
\Omega, & x_i \ge 1,\\
A_i^c, & x_i \in [0,1),\\
\emptyset, &x_i < 0.
\end{cases}
}
It follows that the inverse image $A_X(x)$ lies in $\sF$, and hence $X$ is an $\sF$-measurable random vector. 
\end{exmp}
\end{shaded}
\begin{thm}
Consider a probability space $(\Omega, \sF, P)$, and a finite $n \in \N$. 
A mapping $X: \Omega \to \R^n$ is a random vector if and only if $X_i \triangleq \pi_i \circ X: \Omega \to \R$ are random variables for all $i \in [n]$. 
%In particular, $\sigma(X) = \sigma(X_1, \dots, X_n)$.   
\end{thm}
\begin{proof}
We will first show that $X: \Omega \to \R^n$ implies that $\pi_i \circ X$ is a random variable for any $i \in [n]$. 
For any $i \in [n]$ and $x_i \in \R$, we take $x = (\infty, \dots, x_i, \dots, \infty)$. 
This implies that 
%\EQ{
$\pi_i^{-1}(-\infty, x_i] = \R \times \dots \times (-\infty, x_i] \times \dots \times \R \in \sB(\R^n).$
%} 
Further, defining $A_{X_i}(x_i) \triangleq X_i^{-1}(-\infty, x_i]$, 
we observe from the definition of random vectors that 
\EQN{
\label{eqn:RandomVector:MarginalEvent}
A_X(x) = \cap_{j=1}^nX_j^{-1}(-\infty, x_j]
= X_i^{-1}(-\infty, x_i] 
= A_{X_i}(x_i) \in \sF.  
} 

We will next show that if $X_i: \Omega \to \R$ is a random variable for all $i \in [n]$, 
then $X \triangleq (X_1, \dots, X_n): \Omega \to \R^n$ is a random vector. 
For any $x \in \R^n$, we have $A_{X_i}(x_i) = X_i^{-1}(-\infty,x_i] \in \sF$ for all $i \in [n]$, from the definition of random variables. 
From the closure of event set under countable intersections, 
we have 
\EQN{
\label{eqn:RandomVector:JointEvent}
A_X(x) = \cap_{i=1}^nA_{X_i}(x_i) \in \sF.
} 
\end{proof}
\subsection{Distribution of random vectors}
\begin{defn} %[Distribution of random vectors] 
Consider a probability space $(\Omega, \sF, P)$ and a finite $n \in \N$. 
The \textbf{joint distribution function} of a random vector $X: \Omega \to \R^n$ is defined as the mapping $F_X: \R^n \to [0,1]$ such that 
\EQ{
F_X(x) 
%\equiv F_{X_1, \dots, X_n}(x_1, \dots, x_n) 
\triangleq P(A_X(x)) 
%= P(\set{X_1 \le x_1, \dots, X_n \le x_n}) 
%= P(\cap_{i=1}^nX_i^{-1}(-\infty, x_i])
= P(\cap_{i=1}^nA_{X_i}(x_i)).
}  
\end{defn}
\begin{shaded}
\begin{exmp}[Tuple of indicators] 
Consider a probability space $(\Omega, \sF, P)$, a finite $n \in \N$, and events $A_1, \dots, A_n \in \sF$, 
that define the random vector $X \triangleq (\Ind{A_1}, \dots, \Ind{A_n})$. 
For any $x \in \R^n$, we can define index sets $I_0(x) \triangleq \set{i \in [n]: x_i < 0}$ and $I_1(x) \triangleq \set{i \in [n]: x_i \in [0,1)}$,  
and write the joint distribution function for this random vector $X$ as  
\EQ{
F_X(x) = 
\begin{cases}
1, & I_0(x)\cup I_1(x) = \emptyset,\\
P(\cap_{i \in I_1(x)}A_i^c), & I_0(x) = \emptyset, I_1(x) \neq \emptyset,\\ 
0, & I_0(x) \neq \emptyset.
\end{cases}
}
\end{exmp}
\end{shaded}
\begin{defn}
For a random vector $X: \Omega \to \R^n$ defined on the probability space $(\Omega, \sF, P)$ and $i \in [n]$,  
the distribution of the $i$th random variable $X_i \triangleq \pi_i\circ X: \Omega \to \R$ is called the $i$th \textbf{marginal distribution}, and denoted by $F_{X_i}: \Omega \to [0,1]$. % for all $i \in [n]$. 
\end{defn}
\begin{shaded}
\begin{exmp}[Tuple of indicators] 
Consider a probability space $(\Omega, \sF, P)$, a finite $n \in \N$, and events $A_1, \dots, A_n \in \sF$, 
that define the random vector $X \triangleq (\Ind{A_1}, \dots, \Ind{A_n})$. 
The $i$th marginal distribution is given by  
\EQ{
F_{X_i}(x) = (1-P(A_i))\Ind{[0,1)}(x) + \Ind{[1, \infty)}(x).
}
\end{exmp}
\end{shaded}
\begin{cor}[Marginal distribution] 
Consider a random vector $X: \Omega \to \R^n$ defined on a probability space $(\Omega, \sF, P)$ with the joint distribution $F_X: \R^n \to [0,1]$. 
The $i$th \textbf{marginal distribution} and can be obtained from the joint distribution of $X$ as 
\EQ{
F_{X_i}(x_i) = \lim_{x_j \to \infty, \text{ for all }j \neq i}F_X(x). 
}
\end{cor}
\begin{proof}
For any $i \in [n]$ and $x_i \in \R$,
we have $X_i^{-1}(-\infty, x_i] = A_X(x)$ for $x = (\infty, \dots, x_i, \dots, \infty)$ 
from~\eqref{eqn:RandomVector:MarginalEvent}.  
\end{proof}

\begin{lem}[Properties of the joint distribution function] 
Consider a random vector $X: \Omega \to \R^n$ defined on the probability space $(\Omega,\sF,P)$. 
The associated joint distribution function $F_X: \R^n \to [0,1]$ satisfies the following properties. 
\begin{enumerate}[(i)]
\item For $x, y \in \R^n$ such that $x_i \le y_i$ for each $i \in [n]$, we have $F_X(x) \le F_X(y)$. 
\item The function $F_X(x)$ is right continuous at all points $x \in \R^n$. 
\item The lower limit is $\lim_{x_i \to -\infty}F_X(x) = 0$, and the upper limit is $\lim_{x_i \to \infty, i \in [n]}F_X(x) = 1$.
\end{enumerate}
\end{lem}
\begin{proof}
Consider a random vector $X: \Omega \to \R^n$ defined on the probability space $(\Omega, \sF, P)$ and any $x \in \R^n$. 
\begin{enumerate}[(i)]
\item We can verify that $A_X(x)  = \cap_{i=1}^nA_{X_i}(x_i) \subseteq \cap_{i=1}^nA_{X_i}(y_i) = A(y)$. 
The result follows from the monotonicity of probability measure. 
\item The proof is similar to the proof for single random variable.  
\item The event $A_X(x) = \emptyset$ when $x_i = -\infty$ for some $i \in [n]$ and $A_X(x) = \Omega$ when $x_i = \infty$ for all $i \in [n]$, 
hence the result follow. 
\end{enumerate}
\end{proof}
\begin{shaded}
\begin{exmp}[Probability of rectangular events] 
Consider a probability space $(\Omega, \sF, P)$ and a random vector $X: \Omega \to \R^2$. 
Consider the points $x \le y\in \R^2$ and the events 
\meq{2}{
&B_1 \triangleq \set{x_1 < X_1 \le y_1} = A_{X_1}(y_1)\setminus A_{X_1}(x_1)\in \sF,&
&B_2 \triangleq \set{x_2 < X_2 \le y_2} = A_{X_2}(y_2)\setminus A_{X_2}(x_2)\in \sF.
} 
The marginal probabilities are given by
\eq{
&P(B_1) 
= P(A_{X_1}(y_1))-P(A_{X_1}(x_1)) 
= F_{X_1}(y_1) - F_{X_1}(x_1),\\
&P(B_2) 
= P(A_{X_2}(y_2))-P(A_{X_2}(x_2)) 
= F_{X_2}(y_2) - F_{X_2}(x_2).
}
Writing $x = (x_1,x_2)$ and $y = (y_1,y_2)$, 
we observe that the end points of the rectangular event $B_1\cap B_2$ are points $x, (y_1,x_2), y, (x_1, y_2)$. 
Therefore, we can write this event as 
\EQ{
B_1\cap B_2  = (A_X(y) \setminus A_X(x_1,y_2)) \setminus (A_X(y_1,x_2)\setminus A_X(x)) \in \sF.
}
Hence, we can write the probability of this rectangular event as 
\EQ{
P(B_1\cap B_2) = (F_X(y)-F_X(x_1,y_2)) - (F_X(y_1,x_2)-F_X(x)).
}
\end{exmp}
\end{shaded}

\subsection{Event space generated by random vectors}
\begin{defn} %[Event space generated by random vectors] 
Consider a probability space $(\Omega, \sF, P)$ and a finite $n \in \N$.  
The \textbf{event space generated by a random vector $X:\Omega \to \R^n$} is the smallest $\sigma$-algebra generated by the collection of events $(A_X(x): x \in \R^n)$ and denoted by $\sigma(X) \triangleq \sigma(A_X(x):x \in \R^n)$.  
\end{defn}
\begin{thm}
Consider a probability space $(\Omega, \sF, P)$, a finite $n \in \N$, 
a random vector $X: \Omega \to \R^n$, and its projections $X_i \triangleq \pi_i \circ X$ for all $i \in [n]$. 
Then, $\sigma(X) = \sigma(X_1, \dots, X_n)$.   
\end{thm}
\begin{proof}
Recall that $\sigma(X)$ is generated by the family $(A_X(x): x \in \R^n)$ and $\sigma(X_1, \dots, X_n)$ is generated by the family $(A_{X_i}(x_i): x_i \in \R, i \in [n])$. 
We first show that $A_{X_i}(x_i) = A_X(x)$ for $x = (\infty, \dots, x_i, \dots, \infty)$, and hence $\sigma(X_1, \dots, X_n) \subseteq \sigma(X)$. 
We then show that $A_X(x) = \cap_{i=1}^nA_{X_i}(x_i)$, and hence $\sigma(X) \subseteq \sigma(X_1, \dots, X_n)$. 
\end{proof}
\begin{shaded}
\begin{exmp}[Tuple of indicators] 
Consider a probability space $(\Omega, \sF, P)$, a finite $n \in \N$, and events $A_1, \dots, A_n \in \sF$, 
that define the random vector $X \triangleq (\Ind{A_1}, \dots, \Ind{A_n})$. 
The $\sigma(X) = \sigma(A_i: i \in [n])$. 
\end{exmp}
\end{shaded}
%\begin{rem}
%For a random vector $X: \Omega \to \R^n$ defined on the probability space $(\Omega, \sF, P)$ and any $x \in \R^n$, the event $A_X(x) = \cap_{i=1}^n\set{X_i \le x_i} = \cap_{i=1}^nA_{X_i}(x_i)$. 
%\end{rem}
%%\begin{rem}
%%Hence, the event $A_X(x) = \emptyset$ when $x_i = -\infty$ for some $i \in [n]$ and $A_X(x) = \Omega$ when $x_i = \infty$ for all $i \in [n]$.  
%%That is, 
%%\meq{2}{
%%&\lim_{x_i \to -\infty}F_X(x) = 0, &&\lim_{x_i \to \infty, i \in [n]}F_X(x) = 1.
%%}
%%\end{rem}
\subsection{Independence of random variables}
\begin{defn}
A family of collections of events $(\sA_i \subseteq \sF: i \in I)$ is called independent,
if for any finite set $F \subseteq I$ and $A_i \in \sA_i$ for all $i \in F$, we have 
\EQ{
P(\cap_{i \in F}A_i) = \prod_{i \in F}P(A_i).
} 
\end{defn}
\begin{defn}[Independent and identically distributed]
A random vector $X: \Omega \to \R^n$ defined on the probability space $(\Omega, \sF, P)$ is called \textbf{independent} if 
\EQ{
F_X(x) = \prod_{i=1}^nF_{X_i}(x_i),\text{ for all } x \in \R^n.
}
The random vector $X$ is called \textbf{identically distributed} if each of its components have the identical marginal distribution, i.e.
\EQ{
F_{X_i} = F_{X_1}, \text{ for all } i \in [n]. 
}
\end{defn}
\begin{rem}
Independence of a random vector implies that events $(A_{X_i}(x_i): i \in [n])$ are independent for any $x \in \R^n$. 
Defining families $\sA_i \triangleq (A_{X_i}(x): x \in \R)$ for all $i \in [n]$, 
we observe that the families $(\sA_1, \dots, \sA_n)$ are mutually independent. 
\end{rem}
\begin{rem}
In general, if two collection of events are mutually independent, then the event space generated by them are independent. 
This can be proved using Dynkin's $\pi$-$\lambda$ Theorem. 
\end{rem}
\begin{thm}
For an independent random vector $X: \Omega \to \R^n$ defined on a probability space $(\Omega, \sF, P)$, 
the event spaces generated by its components $(\sigma(X_i): i \in [n])$ are independent.   
%For any Borel measurable sets $A_1,\dots,A_n \in \sB(\R)$, 
%the events $(X_i^{-1}(A_{X_i}) \in \sF: i \in [n])$ are independent. 
\end{thm}
\begin{proof}
For an we define a family of events $\sA_i \triangleq (X_i^{-1}(-\infty, x]: x\in \R)$ for each $i \in [n]$. 
From the definition of independence of random vectors, 
the families $(\sA_i \subseteq \sF: i \in [n])$ are mutually independent. 
Since $\sigma(\sA_i) = \sigma(X_i)$, 
the result follows from the previous remark. 
%that the event spaces $(\sigma(X_i): i \in [n])$ are mutually independent. 
\end{proof}
\begin{defn}[Independent random vectors]
To random vectors $X:\Omega \to \R^n$ and $Y: \Omega\to\R^m$ defined on the same probability space $(\Omega, \sF, P)$ are independent, 
if the collection of events $(A_X(x): x \in \R^n)$ and $(A_Y(y): y \in \R^m)$ are independent, 
where $A_X(x) \triangleq \cap_{i=1}^nX_i^{-1}(-\infty, x_i]$ and $A_Y(y) \triangleq \cap_{i=1}^mY_i^{-1}(-\infty, y_i]$.  
\end{defn}
\begin{shaded}
\begin{exmp}[Independent random vectors] 
Consider a set of vectors $\sX = \set{(0,0,1),(1,0,0)}\subseteq \R^3$. 
Consider two independent coin tosses, such that $\Omega = \set{H,T}^2, \sF = \cP(\Omega)$ and $P(\omega) = p^{k_2(\omega)}(1-p)^{2-k_2(\omega)}$, where $k_2(\omega) = \sum_{i=1}^2\SetIn{\omega_i = H}$. 
We define random vectors 
\begin{xalignat*}{2}
&X = (0,0,1)\SetIn{\omega_1 = H} + (1,0,0)\SetIn{\omega_1 = T},&
&Y = (0,0,1)\SetIn{\omega_2 = H} + (1,0,0)\SetIn{\omega_2 = T}.
\end{xalignat*}
We can verify that $X,Y: \Omega \to \R^3$ are mutually independent random vectors, 
though we can also check that $X_1, X_3$ are dependent random variables and so are $Y_1, Y_3$. 
\end{exmp}
\end{shaded}
\subsection{Discrete random vectors}
\begin{defn}[Discrete random vectors] 
If a random vector $X: \Omega \to \sX_1 \times \dots  \times \sX_n \subseteq \R^n$ 
takes countable values in $\R^n$, 
then it is called a \textbf{discrete random vector}. 
That is, the range of random vector $X$ is countable,
and the random vector is completely specified by the \textbf{probability mass function}
\EQ{
P_X(x) = P(\cap_{i=1}^n\set{X_i = x_i}) \text{ for all }x \in \sX_1\times \dots \times \sX_n.
}
\end{defn}
\begin{rem} 
For an independent discrete random vector $X: \Omega \to \R^n$, we have $P_X(x) = \prod_{i=1}^nP_{X_i}(x_i)$ for each $x \in \R^n$. 
\end{rem}
\begin{shaded}
\begin{exmp}[Multiple coin tosses] 
For a probability space $(\Omega, \sF, P)$, such that $\Omega = \set{H,T}^n, \sF = 2^\Omega, P(\omega) = \frac{1}{2^n}$ for all $\omega \in \Omega$.  

Consider the random vector $X: \Omega \to \R$ such that $X_i(\omega) = \SetIn{\omega_i = H}$ for each $i \in [n]$. 
Observe that $X$ is a bijection from the sample space to the set $\set{0,1}^n$. 
In particular, $X$ is a discrete random variable. 

For any $x \in [0,1]^n$, we can write $N(x) = \sum_{i=1}^n\Ind{[0,1)}(x_i)$. 
Further, we can write the joint distribution as 
\EQ{
F_X(x) = \begin{cases}
1, & x_i \ge 1 \text{ for all }i \in [n],\\
\frac{1}{2^{N(x)}}, & x_i \in [0, 1] \text{ for all } i \in [n],\\
0, & x_i < 0 \text{ for some } i \in [n].
\end{cases}
}
We can derive the marginal distribution for $i$-th component as 
\EQ{
F_{X_i}(x_i) = 
\begin{cases}
1, & x_i \ge 1,\\
\frac{1}{2}, & x_i \in [0, 1),\\
0, & x_i < 0.
\end{cases}
}
Therefore, it follows that $X$ is an \iid random vector. 
\end{exmp}
\end{shaded}

\subsection{Continuous random vectors}
\begin{defn}[Joint density function] 
For jointly continuous random vector $X: \Omega \to \R^n$ with joint distribution function $F_X: \R^n \to [0,1]$, 
there exists a \textbf{joint density function} $f_X: \R^n \to [0,\infty)$ such that $f_X(x) = \frac{d^n}{dx_1\dots dx_n}F_X(x)$, and 
\EQ{
F_X(x) = \int_{u_1\le x_1}du_1\dots \int_{u_n \le x_n} du_n f_X(u_1, \dots, u_n).
}
\end{defn}
\begin{rem}
For an independent continuous random vector $X: \Omega \to \R^n$, we have $f_X(x) = \prod_{i=1}^nf_{X_i}(x_i)$ for all $x \in \R^n$. 
\end{rem}
\begin{shaded}
\begin{exmp}[Gaussian random vectors] 
For a probability space $(\Omega, \sF, P)$, 
Gaussian random vector is a continuous random vector $X: \Omega \to \R^n$ defined by its density function
\EQ{
f_X(x) = \frac{1}{\sqrt{(2\pi)^n\det(\Sigma)}}\exp\left(-\frac{1}{2}(x-\mu)^T\Sigma^{-1}(x-\mu)\right)\text{ for all }x \in \R^n,
}
where the mean vector $\mu \in \R^n$ and the positive definite covariance matrix $\Sigma \in \R^{n \times n}$. 
The components of the Gaussian random vector are Gaussian random variables, 
which are independent when $\Sigma$ is diagonal matrix and are identically distributed when $\Sigma = \sigma^2 I$. 
\end{exmp}
\end{shaded}
%\subsection{Properties of the joint distribution function}

\end{document}
